% Options for packages loaded elsewhere
% Options for packages loaded elsewhere
\PassOptionsToPackage{unicode}{hyperref}
\PassOptionsToPackage{hyphens}{url}
%
\documentclass[
  11pt,
  letterpaper,
  DIV=11,
  numbers=noendperiod,
  twoside]{scrartcl}
\usepackage{xcolor}
\usepackage[left=1.1in, right=1in, top=0.8in, bottom=0.8in,
paperheight=9.5in, paperwidth=7in, includemp=TRUE, marginparwidth=0in,
marginparsep=0in]{geometry}
\usepackage{amsmath,amssymb}
\setcounter{secnumdepth}{3}
\usepackage{iftex}
\ifPDFTeX
  \usepackage[T1]{fontenc}
  \usepackage[utf8]{inputenc}
  \usepackage{textcomp} % provide euro and other symbols
\else % if luatex or xetex
  \usepackage{unicode-math} % this also loads fontspec
  \defaultfontfeatures{Scale=MatchLowercase}
  \defaultfontfeatures[\rmfamily]{Ligatures=TeX,Scale=1}
\fi
\usepackage{lmodern}
\ifPDFTeX\else
  % xetex/luatex font selection
  \setmathfont[]{Garamond-Math}
\fi
% Use upquote if available, for straight quotes in verbatim environments
\IfFileExists{upquote.sty}{\usepackage{upquote}}{}
\IfFileExists{microtype.sty}{% use microtype if available
  \usepackage[]{microtype}
  \UseMicrotypeSet[protrusion]{basicmath} % disable protrusion for tt fonts
}{}
\usepackage{setspace}
% Make \paragraph and \subparagraph free-standing
\makeatletter
\ifx\paragraph\undefined\else
  \let\oldparagraph\paragraph
  \renewcommand{\paragraph}{
    \@ifstar
      \xxxParagraphStar
      \xxxParagraphNoStar
  }
  \newcommand{\xxxParagraphStar}[1]{\oldparagraph*{#1}\mbox{}}
  \newcommand{\xxxParagraphNoStar}[1]{\oldparagraph{#1}\mbox{}}
\fi
\ifx\subparagraph\undefined\else
  \let\oldsubparagraph\subparagraph
  \renewcommand{\subparagraph}{
    \@ifstar
      \xxxSubParagraphStar
      \xxxSubParagraphNoStar
  }
  \newcommand{\xxxSubParagraphStar}[1]{\oldsubparagraph*{#1}\mbox{}}
  \newcommand{\xxxSubParagraphNoStar}[1]{\oldsubparagraph{#1}\mbox{}}
\fi
\makeatother


\usepackage{longtable,booktabs,array}
\newcounter{none} % for unnumbered tables
\usepackage{calc} % for calculating minipage widths
% Correct order of tables after \paragraph or \subparagraph
\usepackage{etoolbox}
\makeatletter
\patchcmd\longtable{\par}{\if@noskipsec\mbox{}\fi\par}{}{}
\makeatother
% Allow footnotes in longtable head/foot
\IfFileExists{footnotehyper.sty}{\usepackage{footnotehyper}}{\usepackage{footnote}}
\makesavenoteenv{longtable}
\usepackage{graphicx}
\makeatletter
\newsavebox\pandoc@box
\newcommand*\pandocbounded[1]{% scales image to fit in text height/width
  \sbox\pandoc@box{#1}%
  \Gscale@div\@tempa{\textheight}{\dimexpr\ht\pandoc@box+\dp\pandoc@box\relax}%
  \Gscale@div\@tempb{\linewidth}{\wd\pandoc@box}%
  \ifdim\@tempb\p@<\@tempa\p@\let\@tempa\@tempb\fi% select the smaller of both
  \ifdim\@tempa\p@<\p@\scalebox{\@tempa}{\usebox\pandoc@box}%
  \else\usebox{\pandoc@box}%
  \fi%
}
% Set default figure placement to htbp
\def\fps@figure{htbp}
\makeatother


% definitions for citeproc citations
\NewDocumentCommand\citeproctext{}{}
\NewDocumentCommand\citeproc{mm}{%
  \begingroup\def\citeproctext{#2}\cite{#1}\endgroup}
\makeatletter
 % allow citations to break across lines
 \let\@cite@ofmt\@firstofone
 % avoid brackets around text for \cite:
 \def\@biblabel#1{}
 \def\@cite#1#2{{#1\if@tempswa , #2\fi}}
\makeatother
\newlength{\cslhangindent}
\setlength{\cslhangindent}{1.5em}
\newlength{\csllabelwidth}
\setlength{\csllabelwidth}{3em}
\newenvironment{CSLReferences}[2] % #1 hanging-indent, #2 entry-spacing
 {\begin{list}{}{%
  \setlength{\itemindent}{0pt}
  \setlength{\leftmargin}{0pt}
  \setlength{\parsep}{0pt}
  % turn on hanging indent if param 1 is 1
  \ifodd #1
   \setlength{\leftmargin}{\cslhangindent}
   \setlength{\itemindent}{-1\cslhangindent}
  \fi
  % set entry spacing
  \setlength{\itemsep}{#2\baselineskip}}}
 {\end{list}}
\usepackage{calc}
\newcommand{\CSLBlock}[1]{\hfill\break\parbox[t]{\linewidth}{\strut\ignorespaces#1\strut}}
\newcommand{\CSLLeftMargin}[1]{\parbox[t]{\csllabelwidth}{\strut#1\strut}}
\newcommand{\CSLRightInline}[1]{\parbox[t]{\linewidth - \csllabelwidth}{\strut#1\strut}}
\newcommand{\CSLIndent}[1]{\hspace{\cslhangindent}#1}



\setlength{\emergencystretch}{3em} % prevent overfull lines

\providecommand{\tightlist}{%
  \setlength{\itemsep}{0pt}\setlength{\parskip}{0pt}}



 


\setlength\heavyrulewidth{0ex}
\setlength\lightrulewidth{0ex}
\usepackage[automark]{scrlayer-scrpage}
\clearpairofpagestyles
\cehead{
  Brian Weatherson
  }
\cohead{
  Logical Positivism and the Rise of Analytic Philosophy
  }
\ohead{\bfseries \pagemark}
\cfoot{}
\makeatletter
\newcommand*\NoIndentAfterEnv[1]{%
  \AfterEndEnvironment{#1}{\par\@afterindentfalse\@afterheading}}
\makeatother
\NoIndentAfterEnv{itemize}
\NoIndentAfterEnv{enumerate}
\NoIndentAfterEnv{description}
\NoIndentAfterEnv{quote}
\NoIndentAfterEnv{equation}
\NoIndentAfterEnv{longtable}
\NoIndentAfterEnv{abstract}
\renewenvironment{abstract}
 {\vspace{-1.25cm}
 \quotation\small\noindent\emph{Abstract}:}
 {\endquotation}
\newfontfamily\tfont{EB Garamond}
\addtokomafont{disposition}{\rmfamily}
\addtokomafont{descriptionlabel}{\rmfamily}
\addtokomafont{title}{\normalfont\itshape}
\let\footnoterule\relax

\makeatletter
\renewcommand{\@maketitle}{%
  \newpage
  \null
  \vskip 2em%
  \begin{center}%
  \let \footnote \thanks
    {\itshape\huge\@title \par}%
    \vskip 0.5em%  % Reduced from default
    {\large
      \lineskip 0.3em%  % Reduced from default 0.5em
      \begin{tabular}[t]{c}%
        \@author
      \end{tabular}\par}%
    \vskip 0.5em%  % Reduced from default
    {\large \@date}%
  \end{center}%
  \par
  }
\makeatother
\RequirePackage{lettrine}

\renewenvironment{abstract}
 {\quotation\small\noindent\emph{Abstract}:}
 {\endquotation\vspace{-0.02cm}}

\setmainfont{EB Garamond Math}[
  BoldFont = {EB Garamond SemiBold},
  ItalicFont = {EB Garamond Italic},
  RawFeature = {+smcp},
]

\newfontfamily\scfont{EB Garamond Regular}[RawFeature=+smcp]
\renewcommand{\textsc}[1]{{\scfont #1}}

\renewcommand{\LettrineTextFont}{\scfont}
% Number theorem-like environments straight through (Theorem 1, Theorem 2, ...)
% rather than by section (1.1, 1.2, ...), which is Quarto's LaTeX default.
% Guarded, since Quarto only defines a counter for environments the document uses.
\makeatletter
\AtBeginDocument{%
  \@for\bw@thm:={theorem,lemma,corollary,proposition,conjecture,definition,%
                 example,exercise,refremark,refsolution}\do{%
    \@ifundefined{c@\bw@thm}{}{\expandafter\counterwithout\expandafter{\bw@thm}{section}}%
  }%
}
\makeatother
\KOMAoption{captions}{tableheading}
\makeatletter
\@ifpackageloaded{caption}{}{\usepackage{caption}}
\AtBeginDocument{%
\ifdefined\contentsname
  \renewcommand*\contentsname{Table of contents}
\else
  \newcommand\contentsname{Table of contents}
\fi
\ifdefined\listfigurename
  \renewcommand*\listfigurename{List of Figures}
\else
  \newcommand\listfigurename{List of Figures}
\fi
\ifdefined\listtablename
  \renewcommand*\listtablename{List of Tables}
\else
  \newcommand\listtablename{List of Tables}
\fi
\ifdefined\figurename
  \renewcommand*\figurename{Figure}
\else
  \newcommand\figurename{Figure}
\fi
\ifdefined\tablename
  \renewcommand*\tablename{Table}
\else
  \newcommand\tablename{Table}
\fi
}
\@ifpackageloaded{float}{}{\usepackage{float}}
\floatstyle{ruled}
\@ifundefined{c@chapter}{\newfloat{codelisting}{h}{lop}}{\newfloat{codelisting}{h}{lop}[chapter]}
\floatname{codelisting}{Listing}
\newcommand*\listoflistings{\listof{codelisting}{List of Listings}}
\makeatother
\makeatletter
\makeatother
\makeatletter
\@ifpackageloaded{caption}{}{\usepackage{caption}}
\@ifpackageloaded{subcaption}{}{\usepackage{subcaption}}
\makeatother
\usepackage{bookmark}
\IfFileExists{xurl.sty}{\usepackage{xurl}}{} % add URL line breaks if available
\urlstyle{same}
% fallback for those not using the hyperref driver hyperxmp:
\makeatletter
\@ifundefined{xmpquote}{\newcommand{\xmpquote}[1]{#1}}{}
\makeatother
\hypersetup{
  pdftitle={Logical Positivism and the Rise of Analytic Philosophy},
  pdfauthor={Brian Weatherson},
  hidelinks,
  pdfcreator={LaTeX via pandoc}}


\title{Logical Positivism and the Rise of Analytic Philosophy}
\author{Brian Weatherson}
\date{2026}
\begin{document}
\maketitle
\begin{abstract}
One of the interesting stories in Sander Verhaegh's \emph{Logical
Positivism: An American History}, is how positivism feeds into the
development of `analytic' philosophy after 1945. A coalition of
`scientific', or later `analytic' philosophers is formed, with
positivists being central to this coalition. What unified this coalition
was not doctrinal agreement, but a commitment to putting logic and
philosophy of science at the centre of philosophy, and to hiring and
promoting people who did that, even if they disagreed on core
philosophical theses. In effect, this became an alliance between
positivists and prominent strands of realism, and within a couple of
decades it took over all the leading philosophy departments in the
United States. I suggest that given the structural advantages these
kinds of coalitions have over philosophical schools with stronger
doctrinal commitments, and the absence of any other coalition until the
rise of `Continental' philosophy many years later, this success was
almost inevitable given the structural arrangements in place in the
1940s. I close with some questions about what would have had to happen
for analytic philosophy not to become the dominant faction, and what
such a world would look like.
\end{abstract}


\setstretch{1.1}
\emph{Logical Positivism} (\citeproc{ref-Verhaegh2026}{Verhaegh 2026})
is a great book.\footnote{In what follows, page references are to this
  book unless otherwise indicated.} It combines careful attention to
archival records, quantitative data, and social context, to tell an
incredibly rich story about the rise, and I guess fall, of logical
positivism in America. I'm going to recommend it to all my students as
crucial background for how philosophy in America got to where it is, and
a correction to several popular narratives.

On the data side, this is one of the best examples you'll see of data
use that highlights simple, understandable statistics. The main
statistic the book uses is simply how often different philosophers are
mentioned in leading American philosophy journals. This is not a simple
statistic to gather; it takes an immense amount of data cleaning and
analysis to tell whether, for example, a given mention is of Roy Wood
Sellars or of Wilfrid Sellars. But Sander Verhaegh's team at Tilburg has
done a mountain of work to get the data in a form where they can tell,
with some confidence, who is being referred to at any point. (They have
also made this data freely available to other researchers, for which
they deserve a lot of credit.) And now, in a book like this, claims
about changes in philosophy can be backed up with these statistics about
how often philosophers are mentioned.

The book also has a number of tables of outliers, e.g., which
philosophers are mentioned most often or, relative to expectation, least
often in a particular period or in a particular journal. These tables
strike me as an absolutely incredible feat. Whenever I build tables of
outliers, they just reveal data errors; the most common words are
usually bits of meta-data that haven't been cleaned from the text. Of
course, it's very easy to understand a table of who was mentioned most
or least often, so it's great for the reader to have something like
this, but I suspect it was an enormous effort to get the tables to be so
useful.

This note will largely be about Verhaegh's discussion of how positivism
fed into the rise of analytic philosophy in America, but first I'll
briefly discuss two smaller points about the book.

\emph{Logical Positivism} is primarily a book about a movement within
academic institutions. But it isn't entirely about that. The effects of
the outside world on positivism are clearly important; we don't get a
generation of exiled empiricists without something causing their exile.
It is discussed well in the book. It isn't just that there were exiles;
who they were, when they came, and where they came from also matter, and
in part 1 of the book we learn a lot more about these questions and why
they matter.

What's more distinctive, though, is the book's discussion of the effects
of positivism on the outside world. In the context of the book,
especially in the discussion of positivism's effect on culture in
chapter 11, positivism is a radical outside force that comes in to shake
things up and is opposed by the conservative mainstream. I had a
question about whether the positivists might have had a lasting and
dramatic effect on American education.

Many American schoolchildren these days are taught that there is a very
sharp fact/opinion distinction. They get worksheets asking them to
classify claims into facts and opinions. These are worksheets, so there
is meant to be a right and wrong answer. Whatever the background of
these worksheets, which I really don't know, they are positivist in
effect. The answers which will be marked correct could have been taken
straight from \emph{Language, Truth and Logic}. That's to say, they have
a very expansive conception of what's `opinion' as opposed to `fact', a
much more expansive one than many of the American positivists Verhaegh
discusses would endorse. The conservative complaints about positivism
that Verhaegh discusses in chapter 11, complaints that positivists were
separating facts and values, would look really weird to modern American
students who have been taught that this distinction is vitally
important. Was getting this into the curriculum a success story for
positivism, or does it have some other history? I'm not sure, but it
feels like an interesting research topic for people who want to look
further into the social effects of positivism.

The other small point I want to note is that it's really pleasing to see
how central \emph{public} universities are to the story in \emph{Logical
Positivism}. The focus is so often on the famous Ivy League schools, and
indeed five of the Ivies are central to Verhaegh's book. But his story
relies on the big public universities too. We see Schlick get a huge
audience at Berkeley very early on in the reception of positivism in
America (144), then Feigl gets hired by Iowa and given more resources at
Minnesota (268, 323). Reichenbach, after many challenges, is hired at
UCLA (147). The central university at the end of the narrative, and
still the home of the most important positivist archive, is Pitt
(395ff). After Charles Stevenson is denied tenure at Yale, he gets a
permanent home at Michigan (253, 268), where I'm based. To this day, our
most prestigious research prize for graduate students is the Stevenson
prize. Stevenson was succeeded at Michigan by Alan Gibbard, so we stayed
a center for non-cognitivist meta-ethics for over half a century.

The main thing I want to talk about is the story Verhaegh tells,
primarily in chapters 9 and 10, of how logical positivism feeds into the
spectacular rise of analytic philosophy in America post-WWII. As with
all plausible narratives of this rise, Verhaegh's starts with the vexed
question of just what this thing, analytic philosophy, is that ended up
winning. The story he gives, and for which he marshals excellent
evidence, is that it becomes in effect the merger of the realist and
positivist traditions. It's a merger of especially the Cambridge realist
tradition with less doctrinaire versions of positivism.

One really striking thing about the narrative is that the label,
`analytic philosophy', gets used for this would-be coalition before
there is any sign of it being a coalition. It's not the only term that
gets used; `scientific' philosophy is just as common. But two projects
by Feigl and Wilfrid Sellars seem to be important in the terminological
shift. One is the collection \emph{Readings in Philosophical Analysis}
(\citeproc{ref-FeiglSellars1949}{Feigl and Sellars 1949}); the other is
the founding of the journal \emph{Philosophical Studies} with the
subtitle \emph{An International Journal for Philosophy in the Analytic
Tradition}. To a large extent, Feigl and Sellars are inventing this
tradition, or perhaps more charitably discovering it, and labeling
it.\footnote{Verhaegh quotes Feigl and Sellars as saying that the two
  traditions ``have increasingly merged'' (323). I think that's true in
  terms of which questions they were asking, which approaches they were
  using, and which answers they took seriously. It wasn't true in terms
  of which answers they gave.} The labeling is a bit of what we'd now
call conceptual engineering. It's massively successful. The forces of
`scientific' philosophy really just run over all before them in the top
American universities.

This coalition, the realists and the positivists, looks so unstable.
What exactly do they have in common? This is why the definitional
question is so hard. The question about commonality can be traced back
to the figures who are conscripted into this new tradition. What exactly
do Carnap, Moore, and Anscombe have in common, which they don't also
share with Peirce?

Here Verhaegh also has a good pair of answers. First, what they have in
common is a connection to \emph{Russell}. Verhaegh shows the positivists
very consciously built a narrative where they were part of a common
tradition with Russell, which helped fend off the charge that they were
just a passing fad (241, 298). Second, there seems to have been
something almost deliberate about writing Peirce into the anti-canon
(334).

Let's step back. One simple narrative, already visible in philosophy's
self-conception, was that there were three primary schools in American
philosophy before positivism arrived. These were idealism, pragmatism,
and realism (\citeproc{ref-Perry1928}{Perry 1928};
\citeproc{ref-Morris1935}{Morris 1935};
\citeproc{ref-Verhaegh2026}{Verhaegh 2026, 154}).\footnote{I'd say that
  is roughly their order of importance at the time, but making that
  claim about their relative position stick would require a major detour
  from this narrative.} There were other smaller schools; in a post-war
survey of the state of philosophy Edgar Sheffield Brightman
(\citeproc{ref-Brightman1947}{1947}) also lists naturalists, Thomists,
personalists and Kantians. But these three were the largest, not least
because they had distinctive sub-schools. It helps to start with these
three.

The coalition the positivists formed with the realists was not
doctrinal; the positivists did not suddenly go around espousing realist
theses.\footnote{For a long time, there was at least a negative thesis
  the members of the coalition had in common: opposition to idealism.
  But as we've seen in recent years, even that can be given up while
  still staying within analytic philosophy, provided one keeps using the
  right logical tools, and responding to the right arguments. See, e.g.,
  Builes (\citeproc{ref-Builes2024}{2024}), and, for a survey of the
  options, Chalmers (\citeproc{ref-Chalmers2019}{2019}).} Rather, it was
administrative. When philosophy departments were expanding in the
post-war years, the realists and positivists worked together to make
sure positions were set aside for `scientific'/`analytic' philosophers.
But at least during the period analytic philosophy was taking over
American philosophy, the positivists were a key part of the coalition.
Why did they enter into \emph{this} coalition, not into a coalition with
one of the other major schools? The short answer is that neither of the
other coalitions worked.

It would have made sense for pragmatists and positivists to work
together. They are both anti-metaphysical. This led people to treat them
similarly at the time. Verhaegh quotes Feigl, Morris, Nagel, and Hook as
treating them as a pair (164, 173, 216). If the two are really
different, where do you put Ramsey? (I come back to him below.) Later in
life, Ayer does an interview with Bryan Magee
(\citeproc{ref-MageeAyer1978}{Magee and Ayer 1978}) where he says
pragmatism was ``very akin to'' positivism.

But this union doesn't work. Verhaegh suggests that it is because the
pragmatists were allergic to dichotomies (185, 187), like the fact/value
one, and that seems plausible. There was also a difference in how they
treated formalism (168). It all seems very contingent; it's easy to
imagine that if Ramsey had lived, he alone might have made the idea of
pragmatist positivism work. But he didn't, and this pairing fails; the
pragmatists at Columbia become some of the staunchest opponents of
positivism.

It's a bit more of a stretch, but you can see how a positivist/idealist
coalition could work. Especially in early versions of positivism, they
both put conscious experience right at the center. Verhaegh quotes
Carnap complaining that people were reading the \emph{Aufbau} as
idealist (151, 155, 210). While Carnap is plausibly right that it's a
misreading, it's not an unintelligible one; there is a synthesis there
if someone wants to develop it. But no one did. It's interesting to
think about what a philosophical landscape would look like where
positivism/idealism vs pragmatism/realism was the main dividing line;
perhaps Ayer vs Austin on sense-data would be taught as the central
debate in philosophy.\footnote{I've elsewhere argued that the standard
  stories of the development of analytic philosophy in the actual world
  don't give enough emphasis to Austin's role
  (\citeproc{ref-Weatherson2007-WEADPW}{Weatherson 2007}). In this
  counterfactual he is much more clearly central.} But that will remain
forever a counterfactual; it's not the way philosophy developed.

Instead, we get a positivist/realist coalition. Doctrinally, this seems
very unstable. But it gets you two advantages, which Verhaegh stresses.
One is that it gets you Russell inside the coalition, without worrying
about Russell's changes of views over time (321). (At least, you don't
have to worry if you ignore Russell's youthful idealism.) More
importantly, it gets you a coalition between philosophers who think that
logic and philosophy of science are central to philosophy. Before WWII,
you already have Morris (\citeproc{ref-Morris1935}{1935}) noting how
much more comfortable the realists are with mathematical approaches to
philosophy than the other schools are. This becomes the defining
characteristic of this new school; not what answers they give to
philosophical questions, but which questions they ask (including
questions from philosophy of science), and how they answer them (with
tools from mathematics and logic). This turned out to be enough.

Here's my contentious claim. Once that was the dividing line, the
victory for the analytics was close to inevitable. We're taught to hate
inevitability claims in history, and rightly so. And I do think it was
really contingent that this became the dividing line. There are many
things that could have happened before WWII that could have upset it. (I
already noted that Ramsey's survival alone might have complicated the
story significantly.) But once those are the sides, I think the
analytics have to win. That's for two reasons.

One is that idealism, especially in America, was dying. And pragmatism
was just not strong enough to fight this new coalition on its own.
Verhaegh has some damning numbers on this. At the key moment, eleven of
Columbia's thirteen faculty were homegrown (187). If Deweyan pragmatism
had really been strong, there would have been graduate students so good
you couldn't help hiring them. There weren't.

The bigger reason is that this coalition left the opposing forces
hopelessly divided. This is already clear in Verhaegh's telling of the
story. At Columbia, the anti-analytics wanted Deweyan pragmatism (260).
At Princeton, they wanted Christian humanism (269). At Yale, they wanted
a department that stayed a stronghold of metaphysics and humanist
ethics. It's not that surprising that it was hard to put together a
unified opposition; if the realists and positivists had to stretch to
find things in common, the idealists, pragmatists, Thomists, etc., who
opposed them had almost nothing in common except a dislike of the new
analytic philosophy.

When analytic philosophy did face a threat, it was from `continental
philosophy'. Now continental philosophy, like analytic philosophy,
covers a multitude of sins. But the advantage each has is that it will
work with a pretty broad range of others under the umbrella, at least
for institution building. That's needed to fight a battle like this, and
the 1950s anti-analytics didn't have it.

Looking at things from a broadly institutional perspective makes it
clear why coalitions, like analytic or continental philosophy, are going
to be well placed to overrun more cohesive schools. For a graduate
student, it is much safer to start working in an analytic program, or a
continental program, than in a program dedicated to pragmatism, or to
Thomism. You can change your mind on key philosophical questions and
still feel at home in the program. If the coalition has a big role in
hiring, which it does, there is a better chance of a good job at the
other end. Both these factors suggest that once a coalition like this
gets rolling, it will start sweeping up graduate students. Indeed,
that's basically what happened both when the `analytic' coalition began
to form after WWII and when the `continental' coalition began to form a
decade or two later.

There is also an advantage at the hiring end. Academics are always fond
of hiring people who share their views, at least about what the
interesting questions are, if not always about what the right answers
are. (Though the latter happens too, of course.) It's sometimes a
challenge, but rarely impossible, to get academics on board with the
idea of hiring in an area they are familiar with but don't think much
of. It's much harder to get them to hire in areas that they are
completely unfamiliar with, which they didn't study in graduate school,
and which they don't interact with. One effect of this is that a school
of philosophy centered around a few universities, even prominent ones,
will lose out if students elsewhere are not learning about it: by the
time those students are involved in hiring, they will have no interest
at all in hiring from that school.

The coalitional approach avoids this problem. The realists and
positivists who were trained by post-war scientific philosophers
disagreed on any number of doctrinal points. But they knew about each
other; indeed, they usually had friends from each side of the various
disputes, and they were willing to hire people who looked like their
teachers or classmates, even if they had disagreed with those teachers
or classmates. (Again, this isn't a point about analytic philosophy; I
think you can tell the same story about continental philosophy, and how
the existence of continental programs creates philosophers who are
willing to hire from the different traditions within continental
philosophy.)

So I'll end with a question, one that I guess relates to the discussion
of explanation in \emph{Logical Positivism}. I think counterfactuals are
crucial to explanation. If we want to explain why analytic philosophy
became dominant, we need to ask what the philosophical world would have
looked like if it never became dominant. My question is, what does that
world look like? And how late can it reasonably and feasibly differ from
actual history? My guess is that the deviation from actual history has
to come relatively early; by 1950, the outcome is baked in, even though
the analytics don't actually take over for several more years. The more
expansive analytic programs are going to have an easier job attracting
good graduate students, and then those students will be more willing to
hire from within this new `analytic tradition'. I'd like to be convinced
I'm wrong; I get nervous with inevitability claims! But this
institutional story seems fairly compelling in retrospect.

As will hopefully be clear by now, \emph{Logical Positivism} is not just
full of fascinating details about how positivism developed in America;
it also tells a compelling story which opens up new ways of looking at
the recent history of philosophy. It's a fantastic book, and I hope it's
really widely read.

\protect\phantomsection\label{refs}
\begin{CSLReferences}{1}{0}
\bibitem[\citeproctext]{ref-Brightman1947}
Brightman, Edgar Sheffield. 1947. {``Philosophy in the {United States}
1939--1945.''} \emph{The Philosophical Review} 56 (4): 390--405. doi:
\href{https://doi.org/10.2307/2181317}{10.2307/2181317}.

\bibitem[\citeproctext]{ref-Builes2024}
Builes, David. 2024. {``Modal Idealism.''} In \emph{Oxford Studies in
Philosophy of Mind}, edited by Uriah Kriegel, 4:46--107. Oxford: Oxford
University Press. doi:
\href{https://doi.org/10.1093/9780198924159.003.0002}{10.1093/9780198924159.003.0002}.

\bibitem[\citeproctext]{ref-Chalmers2019}
Chalmers, David J. 2019. {``Idealism and the Mind-Body Problem.''} In
\emph{The {Routledge} Handbook of Panpsychism}, edited by William
Seager, 353--73. London: Routledge. doi:
\href{https://doi.org/10.4324/9781315717708-28}{10.4324/9781315717708-28}.

\bibitem[\citeproctext]{ref-FeiglSellars1949}
Feigl, Herbert, and Wilfrid Sellars, eds. 1949. \emph{Readings in
Philosophical Analysis}. New York: Appleton-Century-Crofts.

\bibitem[\citeproctext]{ref-MageeAyer1978}
Magee, Bryan, and A. J. Ayer. 1978. {``Logical Positivism and Its
Legacy: {Bryan Magee} Talked to {A. J. Ayer}.''} British Broadcasting
Corporation. \url{https://philpapers.org/rec/MAGLPA-2}.

\bibitem[\citeproctext]{ref-Morris1935}
Morris, Charles W. 1935. {``Some Aspects of Recent American Scientific
Philosophy.''} \emph{Erkenntnis} 5 (1): 142--51. doi:
\href{https://doi.org/10.1007/BF00172290}{10.1007/BF00172290}.

\bibitem[\citeproctext]{ref-Perry1928}
Perry, Ralph Barton. 1928. {``Peace Without Victory---in Philosophy.''}
\emph{Journal of Philosophical Studies} 3 (11): 300--312.
\url{https://www.jstor.org/stable/3745271}.

\bibitem[\citeproctext]{ref-Verhaegh2026}
Verhaegh, Sander. 2026. \emph{Logical Positivism: An American History}.
New York: Oxford University Press.

\bibitem[\citeproctext]{ref-Weatherson2007-WEADPW}
Weatherson, Brian. 2007. {``{Doing Philosophy with Words}.''}
\emph{Philosophical Studies} 135 (3): 429--37. doi:
\href{https://doi.org/10.1007/s11098-006-9026-3}{10.1007/s11098-006-9026-3}.

\end{CSLReferences}



\vspace{12pt} \noindent Draft.


\end{document}
